\section{Related Work}
\label{sec:related}

\paragraph{Volatility forecasting and strong baselines.}
A mature econometric literature sets the bar. GARCH models conditional variance parametrically \citep{bollerslev1986garch,nelson1991egarch}; a parallel tradition treats volatility as observable through realised measures, with Corsi's HAR the parsimonious workhorse \citep{corsi2009har}. These simple models are hard to beat out of sample, and richer predictors add little beyond a strong autoregressive benchmark \citep{hansen2005garch,paye2012dejavol}.

\paragraph{Text in finance and its evaluation gap.}
Closest to us are the designs whose evaluation choices our protocol stress-tests: \citet{kogan2009predicting} regress firm volatility on 10-K word features; FinText forecasts realised volatility with financial word embeddings \citep{rahimikia2021fintext}; FETILDA fine-tunes long-document encoders on 10-K volatility regression \citep{fetilda2022}; and MDRM predicts post-earnings-call volatility from verbal and vocal cues \citep{qin2019mdrm}. None imposes a \emph{recalibrated} price reference, a firm-identity control, or day-clustered inference. The gap is one of protocol, not of baselines: FinText does benchmark against HAR-family models, but earlier designs rarely test significance against the price benchmark.

\paragraph{Neural and LLM representations, and benchmark hygiene.}
Transformer encoders are the strongest general text representations \citep{devlin2019bert,liu2019roberta}, yet their 512-token window falls orders of magnitude short of the median long-form filing; sparse-attention remedies \citep{beltagy2020longformer} are benchmarked almost exclusively on classification rather than on noisy financial regression. Prompted instruction LLMs have recently been proposed as market forecasters \citep{lopezlira2023chatgpt}, raising twin worries of contamination and elicitation fragility \citep{sainz2023contamination}. We add temperature-zero prompting of a 32B model \citep{qwen3_2025} to the spectrum, disciplined by contamination, era-split, paraphrase, repeat-decode, and cross-family checks. Because firms file repeatedly and volatility is firm-persistent, an expressive representation can encode \emph{which firm is speaking} rather than \emph{what is said} \citep{geirhos2020shortcut}. On earnings calls, \citet{yu2025samecompany} establish this directly: same-ticker transcript embeddings are markedly more similar than cross-ticker ones, and their identity-only baseline already matches or beats transcript models. What that leaves open is a calibrated forecast-level verdict (recalibrated reference, day-clustered inference, a stated MDE), which our protocol supplies, and on which MAEC's published gain reprices to nothing detectable. The identity term is a firm fixed effect; our contribution is not that term, nor a new architecture, but the reusable audit protocol that measures its share and travels. To our knowledge, no prior forecast-level evaluation jointly imposes a survivorship-free point-in-time panel, a recalibrated reference interval closed by a maximal price pool, day-clustered inference with Holm and placebo gating, and a model spectrum from dictionaries to a prompted LLM.
